%%%%%%%%%%%%%%%%%%%%%%%%%%%%%%%%%%%%%%%%%%%%%%%%%%%%%%%%%%%
%%% SELECIONE UM DOS MODELOS DE CITAÇÂO, 
%%% REMOVENDO O COMENTÁRIO (%) DA RESPECTIVA LINHA
%%%%%%%%%%%%%%%%%%%%%%%%%%%%%%%%%%%%%%%%%%%%%%%%%%%%%%%%%%%
%\documentclass[tcc]{tcc-poli}   %% citações em formato numérico
\documentclass[tcc,alf]{tcc-poli}  %% citações em formato alfabético

%%%%%%%%%%%%%%%%%%%%%%%%%%%%%%%%%%%%%%%%%%%%%%%%%%%%%%%%%%%
% NÂO ALTERE AS CONFIGURAÇÕES E PACOTES ATÈ O TRECHO DO 
% CÓDIGO/TEXTO INDICADO PARA TAL
%%%%%%%%%%%%%%%%%%%%%%%%%%%%%%%%%%%%%%%%%%%%%%%%%%%%%%%%%%%

%% Pacotes
\usepackage[utf8]{inputenc}
\usepackage{float} 
\usepackage{appendix}
\usepackage{pdflscape}
\usepackage{multirow}
\usepackage{longtable}
\usepackage{graphics}
\usepackage{times}
\usepackage{enumerate}
\usepackage[hidelinks]{hyperref}
\usepackage{epstopdf}
\usepackage{lipsum}
\usepackage{amssymb}
\usepackage{fancyvrb} 
\usepackage{etoolbox}
\usepackage{siunitx}
\usepackage{listings}
\usepackage{caption}
\usepackage{subcaption}
\usepackage[normalem]{ulem}
\useunder{\uline}{\ul}{}

\usepackage{graphicx} %package to manage images
\graphicspath{ {./Figuras/} } %% pasta para organizar as figuras

%%% AJUSTE DOS TAMANHOS E ESTILOS DOS CABECALHOS
%%% NÂO MODIFICAR
\usepackage[big]{titlesec}
\titleformat{\chapter}
{\normalfont\LARGE\bfseries}{\thechapter}{1em}{}
\titleformat{\section}
{\normalfont\Large\bfseries}{\thesection}{1em}{}
\titleformat{\subsection}
{\normalfont\large\bfseries}{\thesubsection}{1em}{}
\titleformat{\subsubsection}
{\normalfont\normalsize\bfseries}{\thesubsubsection}{1em}{}
\titleformat{\paragraph}[runin]
{\normalfont\normalsize\bfseries}{\theparagraph}{1em}{}
\titleformat{\subparagraph}[runin]
{\normalfont\normalsize\bfseries}{\thesubparagraph}{1em}{}
% \titleformat*{\chapter}{\LARGE\bfseries}
% \titleformat*{\section}{\LARGE\bfseries}
% \titleformat*{\subsection}{\Large\bfseries}
% \titleformat*{\subsubsection}{\large\bfseries}
% \titleformat*{\paragraph}{\large\bfseries}
% \titleformat*{\subparagraph}{\large\bfseries}

% para lista de simbolos
\usepackage[intoc]{nomencl}
\makenomenclature

% para lista de abreviaturas
 \usepackage[nomain,acronym,xindy,toc,noredefwarn]{glossaries}
 \makeglossaries

% ABNTEX ...
%% Em próximas versões tentar usar o BIBLATEX em susbstituição do abntex2cite
%\usepackage[
%style = abnt, % Sistema alfabético
% style = abnt-numeric, % Sistema numérico
% style = abnt-ibid, % Notas de referência
%]{biblatex}
%% NÃO alterar os comandos abaixo
\ifx\AbntTexType\StringNum
  \usepackage[num]{abntex2cite}
  \citebrackets[] 
\else
  \usepackage[alf]{abntex2cite}
\fi

%%%%%%%%%%%%%%%%%%%%%%%%%%%%%%%%%%%%%%%%%%%%%%%%%%%%%%%%%%%%
%% A PARTIR DAQUI FAÇA AS ALTERAÇÔES NOS RESPECTIVOS CAMPOS
%%%%%%%%%%%%%%%%%%%%%%%%%%%%%%%%%%%%%%%%%%%%%%%%%%%%%%%%%%%%
%
%% Definicoes
\def\dnl{\mbox{DNL }}
\def\DNL{\mbox{DNL }}

\begin{document}

  %% preencha aqui o Título da Dissertação 
  \title{Título do Trabalho}
  
  %% Título em Ingles
  \foreigntitle{Template for Graduation Project - POLI/UFRJ}
  %% Autor da dissertação
   \author{Nome do(a) Autor(a)} {Sobrenome}
  %% Orientador(es) - adicionar outras linhas caso haja mais orientadores
  \advisor{Prof.}{Nome}{e Sobrenome do(a) Orientador(a)}{D.Sc.} % 
  %\advisor{Prof.}{Nome}{e Sobrenome do(a) Orientador(a) 2}{D.Sc.}  % se houver 2 orientadores
  %\advisor{Prof.}{Nome}{Sobrenome}{D.Sc.}
  
  %% Membros da banca 
  \examiner{Prof.}{Nome completo do membro 1}{D.Sc.}
  \examiner{Prof.}{Nome completo do membro 2}{Ph.D.}
  \examiner{Prof.}{Nome completo do membro 3}{Ph.D.}
  \examiner{Prof.}{Nome completo do membro 4}{D.Sc.}
  % \examiner{Prof.}{Ana Paula Gama}{D.Sc.}

  %% selecionar um Curso (remover comentário do respectivo curso):
  %\engcourse{AMB} % Engenharia Ambiental
  \engcourse{ECV} % Engenharia Civil
  %\engcourse{ECI} % Engenharia de Computação e Informação
  %\engcourse{ECA} % Engenharia de Controle e Automação
  %\engcourse{ELT} % Engenharia Eletrônica e de Computação
  %\engcourse{ELE} % Engenharia Elétrica
  %\engcourse{MAT} % Engeharia de Materiais
  %\engcourse{MET} % Engenharia Metalúrgica
  %\engcourse{MEC} % Engenharia Mecânica
  %\engcourse{NAV} % Engenharia Naval e Oceânica
  %\engcourse{NUC} % Engenharia Nuclear
  %\engcourse{PET} % Engenharia de Petróleo
  %\engcourse{PRO} % Engenharia de Produção
  
  %% Mês e Ano da defesa (não precisa alterar)
  \date{\the\month{}}{\the\year{}}

  %% Palavras-chave em ingles ficam no absract.
  \keyword{Palavra-chave1}
  \keyword{Palavra-chave2}
  \keyword{Palavra-chave3}
  \keyword{Palavra-chave4}

  %%%%%%%%%%%%%%%%%%%%%%%%%%%%%%%%%%%%%%%%%%
  %%  PARTE PRÉ-TEXTUAL. (NÃO MODIFICAR)
  \maketitle
  \frontmatter
  
  \include{dedicatoria} %% OPCIONAL 
  \include{agradecimentos}
  \include{resumo}
  \include{abstract}  
  \tableofcontents 
  \listoffigures
  \listoftables
  \listofquadros
  \include{simbolos} 
  \printnomenclature
  \include{abrev}
  \cleardoublepage
  \printglossary[title=Lista de Siglas e Abreviaturas, type=\acronymtype]
  \mainmatter
  %%%%%%%%%%%%%%%%%%%%%%%%%%%%%%%%%%%%%%%%%%
  %% PARTE TEXTUAL
  
  %% Capítulos 
  \include{capitulo01}
  \include{capitulo02}
  \include{capitulo03}
  \include{capitulo04}
  \include{capitulo05}
  
  \backmatter

  %%%%%%%%%%%%%%%%%%%%%%%%%%%%%%%%%%%%%%%%%%%%
  %% REFERÊNCIAS BIBLIOGRÁFICAS: 
  %% veja como fazer as citações em 
  %% http://mirrors.ctan.org/macros/latex/contrib/abntex2/doc/abntex2cite.pdf
  %%%%%%%%%%%%%%%%%%%%%%%%%%%%%%%%%%%%%%%%%%%%
  %% Padroes ABNT  (NÂO MODIFICAR!)
  \ifx\AbntTexType\StringNum
    \bibliographystyle{abntex2-num}
  \else
    \bibliographystyle{abntex2-alf}
  \fi
  %%%%%%%%%%%%%%%%%%%%%%%%%%%%%%%%%%%%%%%%%%%%

  %% Arquivo com os dados das referências
  \bibliography{referencias}

  %% Apêndices e anexos
  \appendix
  \include{appenA}
  \include{appenB}

\end{document}